\documentclass[aspectratio=169, 11pt, xcolor={table}]{beamer}
% Shared style for the algorithmics course slides (same template as the
% ENSB generative-AI course), with C code listings added.
% Dark blue removed: every slide (cover included) uses the light background
% of the template's content slides, with teal + orange as accents.

\usepackage{amsmath}
\usepackage{amssymb}
\usepackage{booktabs}
\usepackage{colortbl}
\usepackage{tabularx}
\usepackage{tikz}
\usepackage[skins]{tcolorbox}
\usepackage{fontawesome5}
\usepackage{listings}
\usepackage{polyglossia}

\usetikzlibrary{positioning, calc, arrows.meta, shapes.geometric}

%------------------------------------------------
% Languages and fonts (compile from slides/ with XeLaTeX or tectonic)
%------------------------------------------------
\setdefaultlanguage[numerals=maghrib]{arabic}
\setotherlanguage{english}
\usefonttheme{serif}

\newfontfamily\arabicfont[
    Script=Arabic, Scale=0.98, Path=../../fonts/, Extension=.ttf,
    UprightFont=*-Regular, BoldFont=*-Bold,
    ItalicFont=*-Italic, SlantedFont=*-Italic,
    BoldItalicFont=*-BoldItalic, BoldSlantedFont=*-BoldItalic]{Amiri}
\newfontfamily\arabicfontsf[
    Script=Arabic, Scale=0.98, Path=../../fonts/, Extension=.ttf,
    UprightFont=*-Regular, BoldFont=*-Bold,
    ItalicFont=*-Italic, SlantedFont=*-Italic,
    BoldItalicFont=*-BoldItalic, BoldSlantedFont=*-BoldItalic]{Amiri}
% Arabic fallback for \ttfamily (needed by the C listings)
\newfontfamily\arabicfonttt[Script=Arabic, Scale=0.98, Path=../../fonts/,
    Extension=.ttf, UprightFont=*-Regular, BoldFont=*-Bold,
    ItalicFont=*-Italic, SlantedFont=*-Italic,
    BoldItalicFont=*-BoldItalic, BoldSlantedFont=*-BoldItalic]{Amiri}
% Latin monospace font for C code (selected directly, so polyglossia
% does not swap it for the Arabic fallback above)
\newfontfamily\codefont[Scale=0.95, BoldFont=lmmonolt10-bold.otf,
    ItalicFont=lmmono10-italic.otf]{lmmono10-regular.otf}
% Heading font
\newfontfamily\khatt[Script=Arabic, Scale=1.05, Path=../../fonts/,
    BoldFont=ae_AlBattar.ttf]{ae_AlBattar.ttf}

%------------------------------------------------
% Colors (template palette without the navy tones)
%------------------------------------------------
\definecolor{brandTeal}{HTML}{5CC8E6}
\definecolor{brandOrange}{HTML}{FFA27F}
\definecolor{brandLight}{HTML}{F8F9FA}
\definecolor{brandDark}{HTML}{263238}
\definecolor{brandMuted}{HTML}{6C757D}
\colorlet{tealDeep}{brandTeal!60!black}
\definecolor{codeGrey}{HTML}{F1F3F5}

\setbeamercolor{background canvas}{bg=brandLight}
\setbeamercolor{normal text}{fg=brandDark}
\setbeamercolor{structure}{fg=brandTeal}
\setbeamertemplate{navigation symbols}{}

% beamer.cls sets \raggedright globally; under RTL that centers short lines
% or pushes them left, so go back to justified text (flush right).
\rightskip=0pt\relax
% Lists call \raggedright too; make it a no-op for RTL slides.
\AtBeginDocument{\renewcommand{\raggedright}{\rightskip=0pt\relax}}

%------------------------------------------------
% Frame title, footer, lists
%------------------------------------------------
\setbeamertemplate{frametitle}{%
    \vspace{0.35cm}
    \noindent{\khatt\bfseries\Large\color{brandDark}\insertframetitle}
    \par
    \vskip-2pt
    \noindent{\color{brandTeal}\rule{\textwidth}{1.5pt}}
    \vspace{0.05cm}
}

% Footer: page number only
\setbeamertemplate{footline}{%
    \hbox to \paperwidth{\hfill
        {\color{brandMuted}\scriptsize\LR{\insertframenumber\,/\,\inserttotalframenumber}}%
        \hspace{0.6cm}}%
    \vspace{0.25cm}}

\setbeamertemplate{itemize item}{\color{brandTeal}\small\faChevronLeft}
\setbeamertemplate{itemize subitem}{\color{brandMuted}\small\faAngleLeft}
\setbeamertemplate{enumerate item}{\color{tealDeep}\bfseries\insertenumlabel)}

% Tables: every cell framed by a visible border
\newcolumntype{R}{>{\raggedleft\arraybackslash}X}
\newcolumntype{C}{>{\centering\arraybackslash}X}
\newcommand{\ras}[1]{\cellcolor{brandTeal!20}\textbf{#1}}
% Colored text inside a table cell (\leavevmode keeps it on the row baseline)
\newcommand{\cell}[2]{\leavevmode\textcolor{#1}{#2}}
\renewcommand{\arraystretch}{1.35}
\setlength{\arrayrulewidth}{0.6pt}
\arrayrulecolor{brandMuted!70}

%------------------------------------------------
% Boxes (the two boxes of the template's 2nd slide)
%------------------------------------------------
\tcbset{
    brandbox/.style={
        enhanced, arc=4pt, outer arc=4pt, boxrule=1.5pt,
        coltitle=white, coltext=brandDark, fonttitle=\khatt\bfseries,
        toptitle=0.8mm, bottomtitle=0.8mm,
        top=1.2mm, bottom=1.2mm, left=3mm, right=3mm,
        before skip=4pt, after skip=4pt}
}
% Teal box: definitions, objectives, summaries
\newtcolorbox{tbox}[2][]{brandbox, colback=brandTeal!10,
    colframe=brandTeal, title={#2}, #1}
% Orange box: warnings
\newtcolorbox{obox}[2][]{brandbox, colback=brandOrange!10,
    colframe=brandOrange, title={#2}, #1}
% Card (template's 3rd slide, in teal)
\newtcolorbox{card}[2][]{enhanced, colback=brandTeal!6, colframe=brandTeal,
    arc=6pt, boxrule=1.2pt, halign=center, halign title=center,
    coltitle=white, coltext=brandDark, fonttitle=\khatt\bfseries,
    top=1.5mm, bottom=1.5mm, left=2mm, right=2mm, title={#2}, #1}
% Prompt text written to the tool
\newtcolorbox{amr}[1][brandTeal]{enhanced, arc=1mm, boxrule=0pt,
    colback=white, overlay={\fill[#1] (frame.north east) rectangle ([xshift=-2.5pt]frame.south east);},
    left=6pt, right=8pt, top=2pt, bottom=2pt,
    before skip=3pt, after skip=3pt}

%------------------------------------------------
% TikZ (Arabic right-to-left inside nodes)
%------------------------------------------------
\newcommand{\nodecolor}{}
\tikzset{
    every node/.append style={
        execute at begin node={\nodecolor\everypar{\beginR}\beginR},
        execute at end node=\endR},
    lawn/.code={\def\nodecolor{\color{#1}}},
    khatwa/.style={draw=#1, fill=#1!12, very thick, rounded corners=5pt,
        minimum height=0.95cm, text width=1.9cm, align=center,
        font=\khatt\bfseries\small},
    khatwa/.default=brandTeal,
    sahm/.style={-{Stealth[length=2.5mm]}, thick, color=brandMuted},
    rang/.style={draw=#1, fill=#1!12},
    raqm/.style={circle, fill=#1, lawn=white, inner sep=1.5pt,
        font=\small\bfseries},
}

% Horizontal chain of steps, right to left:
% \flow[width]{{a}, {b\\c}, {d}}
\newcommand{\flow}[2][1.9cm]{%
    \begin{tikzpicture}[node distance=5mm]
        \foreach \t [count=\i] in {#2} {
            \ifnum\i=1
                \node[khatwa, text width=#1] (f\i) {\t};
            \else
                \pgfmathtruncatemacro\p{\i-1}
                \node[khatwa, text width=#1, left=of f\p] (f\i) {\t};
                \draw[sahm] (f\p) -- (f\i);
            \fi
        }
    \end{tikzpicture}}

%------------------------------------------------
% C code (always left-to-right, inside a fragile frame)
%------------------------------------------------
\lstdefinestyle{cstyle}{
    language=C, basicstyle=\codefont\footnotesize, columns=fullflexible,
    keepspaces=true, showstringspaces=false, upquote=true,
    keywordstyle=\color{tealDeep}\bfseries,
    commentstyle=\color{red!75!black},
    stringstyle=\color{brandOrange!85!black},
    xleftmargin=16pt,
    numbers=left, numberstyle=\tiny\color{brandMuted}, numbersep=7pt,
    aboveskip=0pt, belowskip=0pt,
    morekeywords={bool, true, false}}
% Inline listings must sit in the english environment, not LTR, otherwise
% numbers such as 3.14 are reversed.
% Numbers in Arabic text: \textenglish{1.75} (not \LR).
% Compact code size for long programs (fixed line height, no gaps)
\newcommand{\codesmall}{\fontsize{8}{8}\selectfont}
\newcommand{\codetiny}{\fontsize{7.4}{7.6}\selectfont}
% Inline C: \code{printf}
\newcommand{\code}[1]{\mbox{\textenglish{\codefont\begingroup\color{tealDeep}#1\endgroup}}}
% Program output (console)
\lstdefinestyle{outstyle}{
    language={}, basicstyle=\codefont\footnotesize\color{brandDark},
    columns=fullflexible, keepspaces=true, showstringspaces=false,
    aboveskip=0pt, belowskip=0pt}

% Listings read from files (no fragile frame needed):
%   \ccode{d1/circle.c}   \out{d1/circle.out}
% One box per listing for the grey background (listings' own per-line
% background leaves thin white gaps between lines)
\newtcolorbox{codebox}{enhanced, colback=codeGrey, colframe=codeGrey,
    boxrule=0pt, arc=0pt, left=4pt, right=4pt, top=1pt, bottom=1pt,
    before skip=0pt, after skip=2pt,
    overlay={\fill[brandMuted!45] (frame.south west) rectangle ([xshift=2.5pt]frame.north west);}}
\newtcolorbox{outbox}{enhanced, colback=brandMuted!20, colframe=brandMuted!20,
    boxrule=0pt, arc=0pt, left=4pt, right=4pt, top=1.5pt, bottom=1.5pt,
    before skip=2pt, after skip=2pt}
\newcommand{\ccode}[2][]{\begin{codebox}\begin{english}\lstinputlisting[style=cstyle,#1]{../code/#2}\end{english}\end{codebox}}
\newcommand{\out}[2][]{\begin{outbox}\begin{english}\lstinputlisting[style=outstyle,#1]{../code/#2}\end{english}\end{outbox}}
% Audience line of the cover
\newcommand{\audience}{طلبة السنة الأولى}

%------------------------------------------------
% Cover (template's 1st slide on the light background)
%------------------------------------------------
\newcommand{\cover}[3]{%
\begin{frame}[plain]
    \begin{tikzpicture}[remember picture, overlay]
        \draw[brandTeal, line width=1.5pt, opacity=0.3]
            (current page.south east) ++(-1,1) circle (4cm);
        \draw[brandTeal, line width=0.8pt, opacity=0.25]
            (current page.north west) ++(2,2) circle (6cm);
        \draw[brandOrange, line width=0.5pt, opacity=0.25]
            (current page.north west) ++(2,2) circle (7cm);
        \node[anchor=north east] at ($(current page.north east)+(-1.0,-0.8)$) {
            \begin{tikzpicture}[scale=0.35]
                \fill[brandTeal] (0,0) -- (0.8,1.2) -- (1.6,0) -- cycle;
                \fill[brandOrange] (0.8,1.2) -- (1.6,2.4) -- (2.4,1.2) -- cycle;
            \end{tikzpicture}
        };
    \end{tikzpicture}

    \vspace{1.6cm}
    \begin{minipage}{0.85\textwidth}
        {\color{tealDeep}\khatt\bfseries #1} \\[0.3cm]
        {\raggedleft\color{brandDark}\Huge\khatt\bfseries #2\par} \vspace{0.3cm}
        {\color{brandMuted}\large #3}
    \end{minipage}

    \vspace{1.2cm}
    \begin{minipage}{0.85\textwidth}
        \color{brandMuted}\small
        \begin{tabular}{@{}ll}
            \textbf{الأستاذ:} & د. بن عبد الرزاق يوسف \\
            \textbf{الفئة:}   & \audience \\
        \end{tabular}
    \end{minipage}
\end{frame}}

\title{البنيات الشرطية}

\begin{document}

\cover{درس الفصل الثاني}
{البنيات الشرطية}
{الشرط البسيط، الشرط المركب، والاختيار المتعدد \textenglish{switch}}

\begin{frame}{أهداف الدرس}
    \begin{tbox}{\faBullseye\quad في نهاية هذا الدرس، تستطيع أن:}
        \begin{itemize}
            \item تكتب عبارة منطقية بعوامل المقارنة والعوامل المنطقية.
            \item تستعمل \code{if} و\code{else} و\code{else if}.
            \item تستعمل \code{switch} للاختيار المتعدد.
        \end{itemize}
    \end{tbox}
\end{frame}

\begin{frame}{لماذا البنيات الشرطية؟}
    \vspace{0.3cm}
    \begin{center}
    \begin{tikzpicture}[
        qarar/.style={diamond, aspect=2.2, draw=brandOrange, fill=brandOrange!12,
            very thick, font=\bfseries\small, inner sep=1pt},
        amal/.style={draw=brandTeal, fill=brandTeal!12, very thick, rounded corners=4pt,
            minimum width=2.4cm, minimum height=0.8cm, font=\small}]
        \node[qarar] (q) at (0,0) {الشرط صحيح؟};
        \node[amal] (y) at (3.2,-1.5) {التعليمات 1};
        \node[amal] (n) at (-3.2,-1.5) {التعليمات 2};
        \draw[sahm] (q.east) -| node[pos=0.25, above, font=\small] {نعم} (y.north);
        \draw[sahm] (q.west) -| node[pos=0.25, above, font=\small] {لا} (n.north);
    \end{tikzpicture}
    \end{center}
    \vspace{0.3cm}
    البرنامج لا ينفذ دائمًا التعليمات نفسها: \textbf{يختار} حسب قيمة الشرط.
\end{frame}

\begin{frame}{عوامل المقارنة}
    \begin{center}
    \begin{tabular}{|c|c|c|}
        \hline
        \ras{المعنى} & \ras{\textenglish{C}} & \ras{مثال} \\ \hline
        يساوي         & \code{==} & \code{x == 0} \\ \hline
        لا يساوي      & \code{!=} & \code{x != 0} \\ \hline
        أصغر، أكبر    & \code{< >} & \code{x > 10} \\ \hline
        أصغر أو يساوي، أكبر أو يساوي & \code{<= >=} & \code{x >= 10} \\ \hline
    \end{tabular}
    \end{center}
    نتيجة المقارنة إما \textbf{صحيح} \textenglish{(1)} أو \textbf{خطأ} \textenglish{(0)}.
\end{frame}

\begin{frame}{العوامل المنطقية}
    \begin{center}
    \begin{tabular}{|c|c|c|}
        \hline
        \ras{العامل} & \ras{\textenglish{C}} & \ras{يكون صحيحًا إذا} \\ \hline
        و     & \code{\&\&} & كان الشرطان صحيحين \\ \hline
        أو    & \code{||}   & كان أحد الشرطين صحيحًا \\ \hline
        نفي   & \code{!}    & كان الشرط خاطئًا \\ \hline
    \end{tabular}
    \end{center}
\end{frame}

\begin{frame}{جدول الحقيقة}
    \begin{center}
    \begin{tabular}{|c|c|c|c|c|}
        \hline
        \ras{\code{A}} & \ras{\code{B}} & \ras{\code{A \&\& B}} & \ras{\code{A || B}} & \ras{\code{!A}} \\ \hline
        0 & 0 & 0 & 0 & 1 \\ \hline
        0 & 1 & 0 & \cell{tealDeep}{\textbf{1}} & 1 \\ \hline
        1 & 0 & 0 & \cell{tealDeep}{\textbf{1}} & 0 \\ \hline
        1 & 1 & \cell{tealDeep}{\textbf{1}} & \cell{tealDeep}{\textbf{1}} & 0 \\ \hline
    \end{tabular}
    \end{center}
\end{frame}

\begin{frame}{مثال على العوامل المنطقية}
    \begin{columns}[T]
        \column{0.5\textwidth}
        \ccode{d2/and.c}
        \column{0.45\textwidth}
        \begin{itemize}
            \item العلامة صالحة إذا كانت بين \textenglish{0} و\textenglish{20}.
            \item الجواب نعم إذا كتب المستخدم \code{y} أو \code{Y}.
        \end{itemize}
    \end{columns}
\end{frame}

\begin{frame}{الشرط البسيط}
    \begin{center}
    \begin{minipage}{0.62\textwidth}
        \ccode{d2/if.c}
    \end{minipage}
    \end{center}
    \vspace{0.3cm}
    إذا كان الشرط صحيحًا تُنفذ التعليمات، وإلا لا يحدث شيء.
\end{frame}

\begin{frame}{الشرط المركب: إذا... وإلا}
    \begin{center}
    \begin{minipage}{0.62\textwidth}
        \ccode{d2/ifelse.c}
    \end{minipage}
    \end{center}
    \vspace{0.3cm}
    تُنفذ \textbf{إحدى} الكتلتين فقط، حسب قيمة الشرط.
\end{frame}

\begin{frame}{تنفيذ: زوجي أم فردي؟}
    \vspace{0.6cm}
    \begin{center}
        \begin{minipage}{0.45\textwidth}
            \out{d2/ifelse.out}
        \end{minipage}
    \end{center}
    \vspace{0.3cm}
    \code{7 \% 2} يساوي \textenglish{1}، إذن الشرط \code{n \% 2 == 0} خاطئ، فيُنفذ \code{else}.
\end{frame}

\begin{frame}{الشروط المتعددة: \textenglish{else if}}
    \begin{center}
    \begin{minipage}{0.62\textwidth}
        \ccode{d2/elseif.c}
    \end{minipage}
    \end{center}
\end{frame}

\begin{frame}{كيف تُقرأ \textenglish{else if}؟}
    \begin{center}
    \begin{tikzpicture}[
        qarar/.style={diamond, aspect=2, draw=brandOrange, fill=brandOrange!15,
            very thick, inner sep=1pt, minimum width=2.9cm, font=\small},
        natija/.style={draw=brandTeal, fill=brandTeal!15, very thick,
            rounded corners=4pt, minimum width=2.3cm, minimum height=0.75cm,
            font=\small\bfseries},
        fail/.style={natija, draw=brandMuted, fill=brandMuted!15},
        lafta/.style={font=\scriptsize\bfseries, lawn=brandMuted}]
        \node[qarar] (q1) at (0,0)     {\code{note >= 16}};
        \node[qarar] (q2) at (-3.6,0)  {\code{note >= 12}};
        \node[qarar] (q3) at (-7.2,0)  {\code{note >= 10}};
        \node[fail]  (r4) at (-10.4,0) {\textenglish{Fail}};
        \node[natija] (r1) at (0,-2)    {\textenglish{Very good}};
        \node[natija] (r2) at (-3.6,-2) {\textenglish{Good}};
        \node[natija] (r3) at (-7.2,-2) {\textenglish{Pass}};
        \node[font=\small\bfseries] (st) at (0,1.5) {\code{note = 13}};
        \draw[sahm] (st) -- (q1);
        \foreach \a/\b in {q1/q2, q2/q3, q3/r4}
            \draw[sahm] (\a) -- node[lafta, above] {لا} (\b);
        \foreach \a/\b in {q1/r1, q2/r2, q3/r3}
            \draw[sahm] (\a) -- node[lafta, right] {نعم} (\b);
        \draw[brandOrange, very thick, rounded corners=6pt, dashed]
            ($(r2.south west)+(-0.15,-0.15)$) rectangle ($(r2.north east)+(0.15,0.15)$);
    \end{tikzpicture}
    \end{center}
    نختبر الشروط بالترتيب، ونتوقف عند \textbf{أول شرط صحيح}.
    مع \code{note = 13} تُعرض \code{Good}.
\end{frame}

\begin{frame}{مثال: أكبر ثلاثة أعداد}
    \begin{columns}[T]
        \column{0.55\textwidth}
        \ccode{d2/max3.c}
        \column{0.4\textwidth}
        \out{d2/max3.out}
    \end{columns}
    \vspace{0.3cm}
    نفترض أن \code{a} هو الأكبر، ثم نقارنه بالعددين الآخرين.
\end{frame}

\begin{frame}{الاختيار المتعدد}
    \begin{tbox}{\faLightbulb\quad متى نستعمل \textenglish{switch}؟}
        عندما نقارن متغيرًا واحدًا بعدة \textbf{قيم ثابتة}،
        من النوع \code{int} أو \code{char}.
    \end{tbox}
\end{frame}

\begin{frame}{بنية \textenglish{switch}}
    \begin{center}
    \begin{minipage}{0.62\textwidth}
        \ccode{d2/switch.c}
    \end{minipage}
    \end{center}
    \vspace{0.3cm}
    \code{default} تُنفذ إذا لم توافق القيمة أي حالة.
\end{frame}

\begin{frame}{دور \textenglish{break}}
    \begin{columns}[T]
        \column{0.5\textwidth}
        \ccode{d2/nobreak.c}
        \column{0.4\textwidth}
        \out{d2/nobreak.out}
    \end{columns}
    \vspace{0.3cm}
    \begin{obox}{\faExclamationTriangle\quad انتبه}
        دون \code{break} يواصل البرنامج تنفيذ الحالات التالية.
    \end{obox}
\end{frame}

\begin{frame}{تجميع عدة حالات}
    \begin{columns}[T]
        \column{0.5\textwidth}
        \ccode{d2/grouped.c}
        \column{0.45\textwidth}
        نكتب عدة حالات متتالية لتنفيذ التعليمات نفسها.
    \end{columns}
\end{frame}

\begin{frame}{مثال: آلة حاسبة}
    \begin{center}\begin{minipage}{0.75\textwidth}\ccode[basicstyle=\codefont\codesmall]{d2/calc.c}\end{minipage}\end{center}
\end{frame}

\begin{frame}{تنفيذ الآلة الحاسبة}
    \vspace{0.6cm}
    \begin{center}
        \begin{minipage}{0.45\textwidth}
            \out{d2/calc.out}
        \end{minipage}
    \end{center}
    \vspace{0.3cm}
    يقرأ البرنامج عددين وعملية، ثم ينفذ الحالة المناسبة.
\end{frame}

\begin{frame}{\textenglish{if} أم \textenglish{switch}؟}
    \begin{center}
    \begin{tabular}{|c|c|c|}
        \hline
        \ras{} & \ras{\code{if}} & \ras{\code{switch}} \\ \hline
        نوع الشرط   & أي عبارة منطقية & مساواة مع قيم ثابتة \\ \hline
        المجالات    & \code{x > 10} ممكن & غير ممكن \\ \hline
        نوع المتغير & كل الأنواع & \code{int} و\code{char} فقط \\ \hline
    \end{tabular}
    \end{center}
\end{frame}

\begin{frame}{أمثلة تطبيقية}
    \begin{tbox}{\faLaptopCode\quad عشرة أمثلة من مجالات مختلفة}
        نطبق ما تعلمناه في مجالات من الحياة اليومية:
        الصحة، الرياضة، الطقس، النقل، التجارة الإلكترونية، ألعاب الفيديو، الأمن السيبراني، الهواتف الذكية، الطاقة الشمسية، السيارات ذاتية القيادة.
    \end{tbox}
\end{frame}

\begin{frame}[t]{المثال الأول: الصحة}
    \begin{columns}[T]
        \column{0.64\textwidth}
        \ccode[basicstyle=\codefont\codesmall]{d2/ex01.c}
        \column{0.33\textwidth}
        {\small\textbf{تصنيف مؤشر كتلة الجسم}: نصنف الحالة الصحية حسب قيمة المؤشر.\par}
        \vspace{0.3cm}
        \out[basicstyle=\codefont\scriptsize\color{brandDark}]{d2/ex01.out}
    \end{columns}
\end{frame}

\begin{frame}[t]{المثال الثاني: الرياضة}
    \begin{columns}[T]
        \column{0.64\textwidth}
        \ccode[basicstyle=\codefont\codesmall]{d2/ex02.c}
        \column{0.33\textwidth}
        {\small\textbf{نتيجة مباراة}: نحسب نقاط الفريق: 3 للفوز، 1 للتعادل، 0 للخسارة.\par}
        \vspace{0.3cm}
        \out[basicstyle=\codefont\scriptsize\color{brandDark}]{d2/ex02.out}
    \end{columns}
\end{frame}

\begin{frame}[t]{المثال الثالث: الطقس}
    \begin{columns}[T]
        \column{0.64\textwidth}
        \ccode[basicstyle=\codefont\codesmall]{d2/ex03.c}
        \column{0.33\textwidth}
        {\small\textbf{إنذار الحرارة}: نختار مستوى الإنذار حسب درجة الحرارة.\par}
        \vspace{0.3cm}
        \out[basicstyle=\codefont\scriptsize\color{brandDark}]{d2/ex03.out}
    \end{columns}
\end{frame}

\begin{frame}[t]{المثال الرابع: النقل}
    \begin{columns}[T]
        \column{0.64\textwidth}
        \ccode[basicstyle=\codefont\codesmall]{d2/ex04.c}
        \column{0.33\textwidth}
        {\small\textbf{رادار السرعة}: نحدد المخالفة حسب تجاوز السرعة المسموح بها.\par}
        \vspace{0.3cm}
        \out[basicstyle=\codefont\scriptsize\color{brandDark}]{d2/ex04.out}
    \end{columns}
\end{frame}

\begin{frame}[t]{المثال الخامس: التجارة الإلكترونية}
    \begin{columns}[T]
        \column{0.64\textwidth}
        \ccode[basicstyle=\codefont\codesmall]{d2/ex05.c}
        \column{0.33\textwidth}
        {\small\textbf{التوصيل المجاني}: التوصيل مجاني إذا تجاوز مبلغ الطلب 5000 دج.\par}
        \vspace{0.3cm}
        \out[basicstyle=\codefont\scriptsize\color{brandDark}]{d2/ex05.out}
    \end{columns}
\end{frame}

\begin{frame}[t]{المثال السادس: ألعاب الفيديو}
    \begin{columns}[T]
        \column{0.64\textwidth}
        \ccode[basicstyle=\codefont\codesmall]{d2/ex06.c}
        \column{0.33\textwidth}
        {\small\textbf{تحريك شخصية}: نحرك الشخصية حسب الزر المضغوط، باستعمال \code{switch}.\par}
        \vspace{0.3cm}
        \out[basicstyle=\codefont\scriptsize\color{brandDark}]{d2/ex06.out}
    \end{columns}
\end{frame}

\begin{frame}[t]{المثال السابع: الأمن السيبراني}
    \begin{columns}[T]
        \column{0.64\textwidth}
        \ccode[basicstyle=\codefont\codesmall]{d2/ex07.c}
        \column{0.33\textwidth}
        {\small\textbf{قوة كلمة المرور}: نقيّم كلمة المرور حسب طولها وعدد أرقامها وحروفها الكبيرة.\par}
        \vspace{0.3cm}
        \out[basicstyle=\codefont\scriptsize\color{brandDark}]{d2/ex07.out}
    \end{columns}
\end{frame}

\begin{frame}[t]{المثال الثامن: الهواتف الذكية}
    \begin{columns}[T]
        \column{0.64\textwidth}
        \ccode[basicstyle=\codefont\codesmall]{d2/ex08.c}
        \column{0.33\textwidth}
        {\small\textbf{وضع البطارية}: نختار وضع تشغيل الهاتف حسب مستوى البطارية.\par}
        \vspace{0.3cm}
        \out[basicstyle=\codefont\scriptsize\color{brandDark}]{d2/ex08.out}
    \end{columns}
\end{frame}

\begin{frame}[t]{المثال التاسع: الطاقة الشمسية}
    \begin{columns}[T]
        \column{0.64\textwidth}
        \ccode[basicstyle=\codefont\codesmall]{d2/ex09.c}
        \column{0.33\textwidth}
        {\small\textbf{الإنتاج حسب الفصل}: نحسب إنتاج لوح شمسي (3 كيلوواط) حسب ساعات الشمس في كل فصل.\par}
        \vspace{0.3cm}
        \out[basicstyle=\codefont\scriptsize\color{brandDark}]{d2/ex09.out}
    \end{columns}
\end{frame}

\begin{frame}[t]{المثال العاشر: السيارات ذاتية القيادة}
    \begin{columns}[T]
        \column{0.64\textwidth}
        \ccode[basicstyle=\codefont\codesmall]{d2/ex10.c}
        \column{0.33\textwidth}
        {\small\textbf{إشارة المرور}: تقرأ السيارة لون الإشارة وتختار التصرف المناسب.\par}
        \vspace{0.3cm}
        \out[basicstyle=\codefont\scriptsize\color{brandDark}]{d2/ex10.out}
    \end{columns}
\end{frame}


\begin{frame}{أخطاء شائعة}
    \begin{obox}{\faExclamationTriangle\quad تجنّب هذه الأخطاء}
        \begin{itemize}
            \item كتابة \code{=} (تعيين) بدل \code{==} (مقارنة).
            \item وضع \code{;} بعد الشرط: \code{if (x > 0);}
            \item نسيان \code{break} في \code{switch}.
        \end{itemize}
    \end{obox}
\end{frame}

\begin{frame}{خلاصة الدرس}
    \begin{tbox}{\faCheckCircle\quad خلاصة الدرس}
        \begin{itemize}
            \item الشرط عبارة منطقية: صحيح أو خطأ.
            \item \code{if / else / else if} لاختيار كتلة من التعليمات.
            \item \code{switch} لمقارنة متغير بعدة قيم ثابتة، مع \code{break}.
        \end{itemize}
    \end{tbox}
\end{frame}

\end{document}
